\subsubsection{Color Threshold Register (COLR)}

Address offset: \textbf{0x24}
\newline
Reset value: 0x0050 0032

\begin{figure}[H]
    \centering
    \begin{bytefield}[bitwidth=\bflenhalfword, endianness=big]{16}
        \bitheader[lsb=16]{16-31} \\
        \bitbox{7}{Reserved} &
        \bitbox{9}{RT} \\
        \bitbox[]{7}{} &
        \bitbox[]{9}{R/W}
    \end{bytefield}
\end{figure}

\begin{figure}[H]
    \centering
    \begin{bytefield}[bitwidth=\bflenhalfword, endianness=big]{16}
        \bitheader{0-15} \\
        \bitbox{7}{Reserved} &
        \bitbox{9}{YT} \\
        \bitbox[]{7}{} &
        \bitbox[]{9}{R/W}
    \end{bytefield}
\end{figure}

\begin{itemize}
    \item \textbf{Bits 31:25 - Reserved}
    \item \textbf{Bits 24:16 - RT: Red Threshold}\\
    Enthält die Grenze in einem 7.2 Festkommaformat in Prozent ab inklusive der die Lautstärke in rot angezeigt wird.\\
    Wird ein Wert größer als 0x64 (100) geschrieben, wird nie eine rote Färbung vorgenommen.\\
    RT hat eine höhere Priorität als YT.
    Sollte also RT kleiner oder gleich YT gewählt werden, wird nie eine gelbe Färbung vorgenommen.
    \item \textbf{Bits 15:9 - Reserved}
    \item \textbf{Bits 8:0 - YT: Yellow Threshold}\\
    Enthält die Grenze in einem 7.2 Festkommaformat in Prozent ab inklusive der die Lautstärke in gelb angezeigt wird.\\
    Wird ein Wert größer als 0x64 (100) geschrieben, wird nie eine gelbe Färbung vorgenommen.\\
    YT hat eine niedrigere Priorität als RT. Also überschreibt RT den Wert YT, wodurch YT größer oder gleich RT zur Folge hat, dass der Balken nie gelb werden kann.

\end{itemize}
